\chapter{Machine Learning para previsão de resultados}
\label{chap:previsao-resultados}

Este capítulo corresponde ao módulo de previsão de resultados de AAUTIA e segue o guia
\path{docs/RESULT_PREDICTION.md}. O objetivo é estimar, antes do início
do jogo, a probabilidade de vitória da equipa da casa. A unidade de
observação é uma partida e o alvo binário indica se a equipa da casa venceu.

\section{Características e informação disponível}

O procedimento utiliza as tabelas de jogos e de estatísticas por equipa
preparadas pelo módulo de preparação de dados. Constrói 46 características, das quais 39 são
diferenças entre médias da equipa da casa e da visitante. Para cada uma
de 13 medidas, calcula médias dos últimos cinco jogos, dos últimos dez
e da temporada em curso. As medidas são vitórias, pontos marcados e
sofridos, margem de pontos, assistências, ressaltos, perdas de bola,
percentagens de lançamentos de campo, triplos e lances livres,
percentagem efetiva de lançamento e eficiências ofensiva e defensiva.

A estimativa de posses de uma equipa é
\[
\widehat{P}=FGA+0{,}44\,FTA-ORB+TOV,
\]
onde \(FGA\) e \(FTA\) são tentativas de lançamentos de campo e lances
livres, \(ORB\) são ressaltos ofensivos e \(TOV\) são perdas de bola.
A média das estimativas das duas equipas define as posses do jogo.
As eficiências correspondem aos pontos marcados ou sofridos por 100
posses estimadas; não são as métricas avançadas oficiais do fornecedor.

As sete características restantes descrevem jogos disputados na época,
dias desde o último jogo, encontros em dias consecutivos e diferença
entre classificações Elo. O descanso é limitado a 14 dias. As médias
reiniciam em cada temporada; o Elo conserva 75\% do afastamento em relação
a 1\,500. A inicialização é 1\,500, com \(K=20\) e vantagem de casa de 65
pontos Elo, parâmetros fixados antes da avaliação.

Os resultados do próprio jogo, identificadores, datas e rótulos das
divisões não entram no modelo. A ausência de histórico permanece como
valor em falta até ao pré-processamento. Este ajusta imputação pela
mediana e indicadores de ausência nos dados de treino, acrescentando
normalização no caso da regressão logística.

\section{Modelos e seleção}

São comparadas três configurações de regressão logística,
\(C\in\{0{,}1;1;10\}\), uma Random Forest com 300 árvores, profundidade
máxima 10 e pelo menos 15 observações por folha, e duas configurações de
\emph{histogram gradient boosting}, com profundidade 3, 200 iterações e
penalização L2 de 1 ou 10. As referências fixas são a frequência histórica
de vitórias em casa e a probabilidade Elo.

O critério de seleção é a menor perda logarítmica na validação:
\[
\mathcal{L}=-\frac{1}{N}\sum_{g=1}^{N}
\left[y_g\log p_g+(1-y_g)\log(1-p_g)\right].
\]
Este critério avalia a qualidade das probabilidades utilizadas pelo
sistema. A exatidão usa o limiar de 0,5; o Brier score mede o erro
quadrático das probabilidades e a ROC-AUC avalia discriminação. Curvas
de calibração e ROC complementam as métricas.

A importância por permutação é calculada na validação, antes do reajuste
final, medindo o aumento de perda após baralhar cada variável.
Características correlacionadas podem partilhar importância; os valores
não representam efeitos causais.

\section{Execução, saídas e limites}

O comando \path{scripts/train_results.py} produz as características,
o dicionário, o modelo guardado, as métricas, previsões de teste e figuras.
O comando \path{scripts/predict_game.py} reutiliza a mesma construção
histórica e rejeita datas dentro do período de treino, equipas desconhecidas
ou identificadores iguais para as duas equipas. Os testes documentados
cobrem alterações de resultados atuais e futuros, exclusão de jogos do
mesmo dia, reinício de temporada e reprodução das probabilidades pelo
modelo guardado.

O modelo não inclui lesões, por falta de validação dos fusos horários dos
relatórios, nem representa explicitamente transferências, cinco inicial
esperado ou recintos neutros. As lacunas das fontes e a bolha de 2019--20
limitam a interpretação da vantagem de casa. Os resultados do
Capítulo~\ref{chap:experiencias} descrevem previsões diárias retrospetivas
numa temporada de teste.
